\section{Deletion lines and concentration of core covers}
\label{sec:deletion}

The deterministic construction asks for rigidity of line covers of the
core $R_0$. We establish that rigidity in two steps. First, after a fixed
number of lines in each direction are deleted, with probability tending
to one every cover of the surviving set by at most $q$ other lines also
covers almost all of the plane. Second, over a prime field, the Sz\H{o}nyi--Weiner stability
theorem puts almost all lines of such an almost-cover in one pencil.

The deletion law matters. Our eventual split lines come from a constrained
selection, so their law is not uniform. We prove the first step for any
law whose probability at each deletion set is at most $\exp(Kq)$ times
its uniform probability, with $K$ fixed. A qualitative estimate for
uniform deletion would not suffice under this change of measure.

Throughout this section, $q$ tends to infinity through prime powers and
$A=\mathbb F_q^2$.  Let $\mathcal L$ be its $q(q+1)$ affine lines.  Primality
will be required only for the final application of projective stability.
Fix an integer $k\ge1$.  For $q\ge k$, let $\mathsf U_{q,k}$ be the uniform
law on sets $D\subseteq\mathcal L$ containing exactly $k$ lines in each
direction.  Define
\[
 R_D=A\setminus\bigcup_{\ell\in D}\ell,
 \qquad i_D(x)=|\{\ell\in D:x\in\ell\}|.
\]
For probability measures on a countable set, total variation distance is
one half of the sum of the absolute coordinate differences.

\begin{theorem}[Robust deletion estimates]\label{thm:deletion}
Fix $k\ge1$ and $K\ge0$.  For each $q$, let $\mathsf Q_q$ be a probability
law on the same sets as $\mathsf U_{q,k}$ such that
\[
 \mathsf Q_q(D)\le e^{Kq}\mathsf U_{q,k}(D)
 \quad\text{for every }D.
\]
Then the following assertions hold as $q\to\infty$.
\begin{enumerate}[label=\textup{(\roman*)}]
\item The empirical probability measure
$q^{-2}\sum_{x\in A}\delta_{i_D(x)}$ converges in probability, in total
variation, to the Poisson distribution with mean $k$.
\item Write $\mu=e^{-k}$.  For every fixed $\lambda,\eta>0$, the probability
that at least $\eta q$ lines $\ell\in\mathcal L\setminus D$ satisfy
\[
 \left|\frac{|\ell\cap R_D|}{q}-\mu\right|\ge\lambda
\]
tends to zero.
\item For every fixed $\delta>0$, the probability that there is a set
$C\subseteq\mathcal L\setminus D$ with $|C|\le q$ which covers $R_D$ and
leaves at least $\delta q^2$ points of $A$ uncovered tends to zero.
\end{enumerate}
\end{theorem}

Part~\textup{(iii)} is the cover statement needed for the construction.
The first two parts provide its input: the surviving density and the
occupancy of all but a sublinear number of lines. We prove them in that
order using one counting principle.

For a proposed witness $C$ of size $m$, expose the lines of $D$ and $C$
in a carefully chosen random order. For a positive proportion of orders,
their point-multiplicity profiles will force more than $m$ exposures at
which the next line has atypical occupancy. In the counting experiment,
where $D$ and $C$ are themselves sampled, each such step has conditional
probability $O(q^{-1})$.
The first $m$ savings pay for the at most $q^m\exp(O(q))$ witnesses;
the remaining positive proportion of steps pays for $e^{Kq}$.
Erd\H{o}s and Lov\'asz~\cite[Section~4]{ErdosLovasz1975} count small
covers by their collinear subsets in a random-line construction;
Kahn~\cite{Kahn1992} subsequently strengthened the random-line result.
The profile count below provides the uniform change-of-measure control
used here.

\subsection{Profiles and incidence identities}

Consider disjoint sets $D,C$ with $D$ as above and $m=|C|\le q$.  Put
$j_C(x)=|\{\ell\in C:x\in\ell\}|$ and define the plane profile and the
profile of a line $\ell\in D\cup C$ by
\[
 P^{(q)}_{ij}=q^{-2}|\{x\in A:i_D(x)=i,\ j_C(x)=j\}|,
 \qquad
 X^\ell_{ij}=q^{-1}|\{x\in\ell:i_D(x)=i,\ j_C(x)=j\}|.
\]
All indices in these arrays are nonnegative integers.  A coordinate with
a negative index is understood to be zero.

The mean of $i+j$ in the plane is $(k(q+1)+m)/q$.  On any one of the
listed lines it is at most
\[
 1+\frac{k(q+1)+m-1}{q},
\]
because every other distinct line meets it in at most one point.  Thus
all these arrays belong to the space $\mathscr P$ of probability arrays
on $\mathbb Z_{\ge0}^2$ whose mean of $i+j$ is at most $M=2k+4$.
Give $\mathscr P$ the topology of coordinatewise convergence.

This space is compact and metrizable.  Indeed the mass outside
$\{i+j\le B\}$ is at most $M/(B+1)$, uniformly over the space.  A diagonal
subsequence of any sequence therefore has total mass one, and its mean
bound follows from Fatou's lemma.  The same tail estimate shows that
coordinatewise convergence in $\mathscr P$ implies convergence in total
variation.  In particular every bounded sum of selected coordinates,
such as $x\mapsto\sum_jx_{0j}$, is continuous.

From any sequence of pairs $(D,C)$, we may consequently extract a
subsequence such that
\begin{equation}\label{eq:profile-limits}
 \frac mq\longrightarrow c,\quad P^{(q)}\longrightarrow P,\quad
 \nu_D^{(q)}:=\frac1q\sum_{\ell\in D}\delta_{X^\ell}
       \Longrightarrow\nu_D,\quad
 \nu_C^{(q)}:=\frac1q\sum_{\ell\in C}\delta_{X^\ell}
       \Longrightarrow\nu_C.
\end{equation}
The last two limits are weak limits of finite measures on $\mathscr P$.
Their masses are $k$ and $c$.  The measure $\nu_D$ is supported on profiles
with $x_{0j}=0$, and $\nu_C$ on profiles with $x_{i0}=0$; these restrictions
are closed.  For each fixed $i,j$, counting incidences before taking the
limit gives
\begin{equation}\label{eq:profile-incidences}
 \int x_{ij}\,d\nu_D(x)=iP_{ij},\qquad
 \int x_{ij}\,d\nu_C(x)=jP_{ij}.
\end{equation}
Summing these identities by Tonelli's theorem yields
\begin{equation}\label{eq:profile-means}
 \sum_{i,j}iP_{ij}=k,\qquad \sum_{i,j}jP_{ij}=c.
\end{equation}
This argument is necessary: convergence in total variation alone would
not justify passage to these unbounded first moments.

The next distinction identifies exactly which line profiles can create
an atypical exposure. If a line of type $D$ sampled the surrounding plane
without bias, its own contribution would merely add one to the first
multiplicity coordinate. The analogous shift for a $C$-line is in the
second coordinate.

Relative to the limiting plane profile $P$, call a profile of type $D$
\emph{invisible} if $x_{ij}=P_{i-1,j}$ for every $i,j$, and a profile of
type $C$ invisible if $x_{ij}=P_{i,j-1}$ for every $i,j$.  A profile which
is not invisible is called visible.  The terminology refers only to the
following exposure test.

\begin{lemma}[Counting by profile detection]\label{lem:profile-counting}
For each $q$, let $\mathscr A_q$ be a deterministic class of pairs $(D,C)$
as above.  Assume there is $a>0$ such that every pair in every class has
either $m=0$ or $m\ge aq$.  Suppose every subsequential limit
\eqref{eq:profile-limits}, along objects in the classes with $q\to\infty$,
satisfies
\begin{equation}\label{eq:visible-excess}
 \nu_D(\text{$D$-visible profiles})+
 \nu_C(\text{$C$-visible profiles})>c.
\end{equation}
There is $b>0$ such that, for all sufficiently large $q$,
\[
 \mathsf U_{q,k}\{D:\text{some }C\text{ has }(D,C)\in\mathscr A_q\}
 \le \exp(-bq\log q).
\]
The constants may depend on $k,a$ and the sequence of classes.
\end{lemma}

\begin{proof}
We first keep $(D,C)$ fixed and show that a positive proportion of
its exposure orders produce more than $m$ marked steps. We then
sample $D$ and $C$ independently and show that so many marks are
rare. Comparing the two experiments will pay for all possible
witness sets $C$.

\paragraph{Exposure of a fixed pair.}
First fix one deterministic object $(D,C)$ and randomize only its
exposure order.  Order $D$ and $C$ independently and uniformly, and set
$n_D=k(q+1)$ and $\gamma=\sqrt2$.  With equal exposure rates, limiting
point-survival probabilities would depend only on $i+j$.
The irrational exponent makes
$i+\gamma j$ distinct for distinct pairs $(i,j)$, allowing the exposure
test to distinguish both coordinates.  Assign the successive $D$-steps the
parameters $d/n_D$ for $d=n_D,n_D-1,\ldots,1$, and the successive
$C$-steps the parameters $(a'/m)^{1/\gamma}$ for
$a'=m,m-1,\ldots,1$.  Omit the second list if $m=0$.  Interleave the lists
in decreasing parameter order, breaking ties deterministically.  At a
step with parameter $z$, the fractions of lines of the two types still
unexposed, including the current line of its type, differ from $z$ and
$z^\gamma$ by at most $O(1/n_D)$ and $O(1/m)$, respectively.

Let $U$ be the points on none of the previously exposed lines.  For
fixed thresholds $\eta_0,\epsilon>0$, mark a step if $z>\eta_0$ and its
line $\ell$ satisfies
\begin{equation}\label{eq:mark}
 \left|\frac{|\ell\cap U|}{q}-\frac{|U|}{q^2}\right|>\epsilon.
\end{equation}
Write $N$ for the number of marks.  We relate its expectation to the
profiles and then count marked exposure sequences.

\paragraph{Detecting the visible profiles.}
For a limit \eqref{eq:profile-limits}, put
\[
 u(z)=\sum_{i,j}P_{ij}z^{i+\gamma j},\qquad
 f_D(x,z)=\sum_{i\ge1,j\ge0}x_{ij}z^{i-1+\gamma j},\qquad
 f_C(x,z)=\sum_{i\ge0,j\ge1}x_{ij}z^{i+\gamma(j-1)}.
\]
These functions are continuous jointly in the displayed variables for
$0<z\le1$ on the relevant profile spaces.  Conditional on a particular
line $\ell$ occupying a specified step of its type, the expectation of
$|U|/q^2$ is $u(z)+o(1)$ and that of $|\ell\cap U|/q$ is
$f_T(X^\ell,z)+o(1)$, uniformly for $z>\eta_0$ and for all target lines of
type $T$.  To verify this, at a $D$-step with $d$ remaining $D$-lines and
$a'$ remaining $C$-lines, a point on $\ell$ of multiplicities $i,j$ is
uncovered with probability
\[
 \frac{(d-1)_{i-1}}{(n_D-1)_{i-1}}
 \frac{(a')_j}{(m)_j}.
\]
Here $(t)_r=t(t-1)\cdots(t-r+1)$, and the second factor is one if $j=0$;
when $m=0$ only $j=0$ occurs.  Off $\ell$, the first factor is
$(d-1)_i/(n_D-1)_i$.  The analogous formulas at $C$-steps shift $j$
instead.  For bounded $i+j$ these expressions have the claimed uniform
power approximations.  The uniform profile mean bound controls the
remaining tails, and the target line itself occupies only $1/q$ of the
plane.  We may retain only a subsequence with $m=0$, or one with $m\ge aq$,
so every denominator in the relevant bounded-multiplicity approximation
tends to infinity.

Both normalized uncovered counts have conditional variance $O(1/q)$.
Indeed the previously exposed lines form uniform subsets of specified
sizes with the target excluded.  Exchanging a selected and an unselected
line changes either normalized count by at most $2/q$: on the target
each other line meets at most one point, and in the plane each line has
$q$ points.  The transposition coupling for sampling without replacement
gives Doob martingale increments at most $2/q$ in $O(q)$ exposures, proving
the variance estimate.  This establishes concentration conditional on
each target, uniformly in the target and its step.

For a visible profile $x$ of type $T$, the functions $f_T(x,z)$ and $u(z)$
agree only on a set of Lebesgue measure zero in $(0,1)$.  For proof, put
$z=e^{-s}$.  Their difference is an analytic generalized Dirichlet series
on $\operatorname{Re}s>0$, with absolutely summable coefficients and
exponents $i+\gamma j$.  These exponents are distinct and locally finite.
If the difference vanishes identically and some coefficient is nonzero,
choose its least exponent, divide by the corresponding exponential, and
let $s\to+\infty$.  Dominated convergence gives a nonzero limit, a
contradiction.  Thus identity of the functions is exactly invisibility.
The zero set of a nonzero analytic function on the real interval is
discrete there and hence null.

At a fixed step the target is uniform among lines of its type.  The
uniform parameter grids converge to the probability densities $1$ for
$D$ and $\gamma z^{\gamma-1}$ for $C$.  Weak convergence of the products
of these grids and the empirical profile measures, followed by the
open-set lower bound and the conditional concentration just proved,
gives
\begin{align}
 \liminf\frac{\mathbb E N}{q}\ \ge\ &
 \iint\mathbf1_{\{\eta_0<z<1,\ |f_D(x,z)-u(z)|>3\epsilon\}}
            \,d\nu_D(x)\,dz \notag\\
 &+\iint\mathbf1_{\{\eta_0<z<1,\ |f_C(x,z)-u(z)|>3\epsilon\}}
            \,d\nu_C(x)\,\gamma z^{\gamma-1}\,dz.
 \label{eq:marks-limit}
\end{align}
The right side increases to the visible mass sum when the thresholds
decrease to zero.

The limiting test now has a uniform consequence. The strict excess in
\eqref{eq:visible-excess} need not have been quantified in the lemma's
hypothesis; compactness supplies the margin needed for counting.
There are fixed $\eta_0,\epsilon,\xi>0$ such that every sufficiently large
object in the classes satisfies
\begin{equation}\label{eq:marks-margin}
 \mathbb E N\ge m+\xi q.
\end{equation}
Otherwise take both thresholds and the proposed margin equal to $1/t$,
choose violating objects with $q_t\to\infty$, and pass to a subsequential
profile limit.  For any fixed positive thresholds, their mark count is
eventually no larger than the count at thresholds $1/t$.  Its limiting
upper expectation divided by $q_t$ is therefore at most $c$.  Equation
\eqref{eq:marks-limit}, followed by decreasing the fixed thresholds to
zero, contradicts \eqref{eq:visible-excess}.  This proves the uniform
margin without assuming a quantitative excess in advance.  Since
$N\le n_D+m$, \eqref{eq:marks-margin} implies
\begin{equation}\label{eq:marks-positive-probability}
 \mathbb P\{N\ge m+\xi q/2\}\ge\frac{\xi q}{2n_D}.
\end{equation}
The probability in this display is only over orderings of the fixed
object, and its lower bound is a positive constant for large $q$.

\paragraph{Counting the marked exposures.}
We have proved that every fixed forbidden pair has a positive probability
of producing more than $m$ marks. We now use a different experiment to
count how often any pair can produce that many marks.  Sample $D$ from
$\mathsf U_{q,k}$, sample $C$ independently as a uniform $m$-subset of
$\mathcal L$, and use the same random orderings.  This experiment is
\emph{not} conditioned on $D\cap C=\varnothing$.  Disjointness will enter
only when we apply \eqref{eq:marks-positive-probability} to an object in
$\mathscr A_q$.

At a $D$-step with $z>\eta_0$, conditional on all past exposures, the
probability of any particular next line is $O(q^{-2})$.  To see this,
sample a uniform ordering of direction tokens, with $k$ copies of each
direction, and choose its offsets uniformly without replacement when a
token recurs.  A direction has conditional probability at most
$k/(\eta_0n_D)$, and any line in that direction then has probability at
most $1/(q-k)$ for large $q$.  Past $C$-exposures do not change this law.
At a $C$-step the same bound follows from uniform sampling without
replacement from $q(q+1)$ lines.

At this point $U$ is fixed by the exposed history.
Lemma~\ref{lem:affine-second-moment} gives, for this arbitrary set,
\begin{equation}\label{eq:affine-variance}
 \sum_{\ell\in\mathcal L}
       \left(|\ell\cap U|-\frac{|U|}{q}\right)^2
       =q|U|-\frac{|U|^2}{q}\le\frac{q^3}{4}.
\end{equation}
It follows that at most $q/(4\epsilon^2)$ lines can cause a mark.  Thus every step has
conditional mark probability at most $B/q$, for a fixed $B$; steps below
the parameter threshold have mark probability zero.  Multiplying these
conditional bounds over any specified collection of steps and taking a
subset union bound over $O(q)$ steps gives
\begin{equation}\label{eq:many-marks-upper}
 \mathbb P\{N\ge m+\xi q/2\}
 \le \exp(O(q))q^{-m-\xi q/2}.
\end{equation}
This uses no independence between marks.

\paragraph{Paying for all witnesses.}
By \eqref{eq:marks-positive-probability}, the probability that the sampled
pair belongs to $\mathscr A_q$ is at most a constant times the right side
of \eqref{eq:many-marks-upper}.  Passing from a random $C$ to the existence
of a witness $C$ costs at most
\[
 \binom{q(q+1)}m\le q^m\exp(O(q)),\qquad 0\le m\le q.
\]
For $m>0$, this follows from $\binom nm\le(en/m)^m$ and
$m\log(q/m)=O(q)$; $m=0$ is immediate.  The $q^m$ factors cancel.
Summing over the at most $q+1$ possible values of $m$ leaves
$\exp(-\xi q\log q/2+O(q))$, which proves the lemma.
\end{proof}

\subsection{The three deletion estimates}

\begin{proof}[Proof of Theorem~\ref{thm:deletion}]
Multiplication of the bound in Lemma~\ref{lem:profile-counting} by $e^{Kq}$
still gives a probability tending to zero.  We will use this observation
for three classes of witnesses.

\paragraph{The multiplicity law.}
For part~\textup{(i)}, fix a positive total variation tolerance and take
$m=0$ and the class where the $D$-multiplicity profile violates that
tolerance from Poisson$(k)$.  In a subsequential limit $j=0$ throughout.
If all $D$-mass is invisible, \eqref{eq:profile-incidences} gives
\[
 iP_{i0}=kP_{i-1,0}\quad(i\ge1).
\]
Normalization then forces $P_{i0}=e^{-k}k^i/i!$, contrary to the fixed
tolerance, which persists by total variation convergence.  Thus some
positive $D$-mass is visible, and the counting lemma applies with $c=0$.
This proves part~\textup{(i)}.

Fix now the particular sequence of dominated laws in the theorem.  By
part~\textup{(i)}, we may choose a deterministic sequence $\epsilon_q$
decreasing to zero sufficiently slowly that, with probability tending
to one, the empirical $D$-multiplicity law lies within $\epsilon_q$ of
Poisson$(k)$.  Such a sequence follows by choosing successive thresholds
$1/t$ after their failure probabilities have become small.  These
restrictions define deterministic classes of outcomes; they do not
condition the independent counting experiment in the lemma.

\paragraph{Line occupancy.}
For part~\textup{(ii)}, it suffices to take $0<\eta<1$.  If at least
$\eta q$ non-$D$ lines have the stated deviation, at least
$m=\lfloor\eta q/2\rfloor$ have deviations of one common sign.  Treat
the two signs separately, and let $C$ be any such $m$-set.  Include the
Poisson-closeness restriction just chosen in the class of pairs.
In every subsequential limit, the $i$-marginal of $P$ is Poisson$(k)$, so
\[
 \sum_jP_{0j}=\mu.
\]
Every limiting $C$-profile has $\sum_jx_{0j}\ge\mu+\lambda$, or every
one has $\sum_jx_{0j}\le\mu-\lambda$, according to the sign.  These are
closed conditions on the compact profile space.  Consequently all
$C$-mass is visible: an invisible $C$-profile has this row sum equal to
$\sum_jP_{0j}=\mu$.

If all $D$-mass were invisible, the incidence identities would give
\begin{equation}\label{eq:poisson-factorization}
 P_{ij}=\frac{k^i}{i!}P_{0j}.
\end{equation}
By Tonelli's theorem and \eqref{eq:profile-means}, this would imply
\[
 \int\sum_jx_{0j}\,d\nu_C(x)
 =\sum_jjP_{0j}
 =e^{-k}\sum_{i,j}jP_{ij}=\mu c.
\]
But the common sign deviation bounds this integral above by
$(\mu-\lambda)c$ or below by $(\mu+\lambda)c$, and
$c=\eta/2>0$.  This is a contradiction.  There is therefore also positive
visible $D$-mass, giving the strict excess required by the counting
lemma.  Its bound, the domination factor, and the vanishing failure
probability of the Poisson restriction prove part~\textup{(ii)}.

\paragraph{A core cover almost covers the plane.}
We have obtained the density and line occupancy needed to analyze a
cover. The occupancy estimates and the budget of $q$ lines force the
$i=0$ row of its limiting profile to be supported on $j=1$. If the cover
still misses a positive proportion of the plane, additional deletion
profiles must then be visible.

A second diagonal choice, using parts~\textup{(i)} and \textup{(ii)},
provides deterministic thresholds tending to zero such that, with
probability tending to one, the Poisson restriction holds and all but
$o(q)$ non-$D$ lines have $|\ell\cap R_D|/q=\mu+o(1)$, with the errors
uniform among the nonexceptional lines.  More explicitly, one may choose
$t=t(q)\to\infty$ slowly so that the number of lines with deviation at
least $1/t$ is smaller than $q/t$, except with probability tending to zero.

Fix $\delta>0$ and consider witnesses for failure of part~\textup{(iii)}
subject to these deterministic restrictions.  Covering $R_D$ by $m$ lines
implies
\[
 \frac mq\ge\frac{|R_D|}{q^2}=\mu+o(1),
\]
so the lower bound on positive $m/q$ required by the counting lemma holds.
In a limit, almost every $C$-profile satisfies
$\sum_jx_{0j}=\mu$, because at most $o(q)$ members of $C$ can be among the
exceptional lines.  Covering gives $P_{00}=0$, and $c\le1$.  Hence
\[
 c\mu
 =\int\sum_jx_{0j}\,d\nu_C
 =\sum_jjP_{0j}
 \ge\sum_jP_{0j}=\mu.
\]
It follows that $c=1$, $P_{01}=\mu$, and $P_{0j}=0$ for $j\ne1$.
Moreover \eqref{eq:profile-incidences} gives
$\int x_{01}\,d\nu_C=P_{01}=\mu$.  Since $x_{01}\le\sum_jx_{0j}=\mu$
almost everywhere and $\nu_C$ has mass one, $x_{01}=\mu$ almost everywhere.
Every such profile is visible, because an invisible $C$-profile would
have $x_{01}=P_{00}=0$.

All $D$-mass cannot be invisible either.  Otherwise
\eqref{eq:poisson-factorization} and the preceding description of the
$i=0$ row would force $P$ to be supported on $j=1$.  This contradicts
$\sum_iP_{i0}\ge\delta$, the limiting fraction of points uncovered by
$C$.  Thus the visible mass is strictly greater than $c$ again.  The
counting lemma applies, and adding back the vanishing probabilities of
the deterministic restrictions proves part~\textup{(iii)}.
\end{proof}

\subsection{From almost covering the plane to a pencil}

The deletion theorem leaves only a geometric question: how can at most
$q$ lines cover all but a small proportion of an affine plane?
An earlier stability lemma of Erd\H{o}s and
Lov\'asz~\cite[Section~4, proof of Theorem~10, part~II]{ErdosLovasz1975}
forces almost all of a set of $q$ projective points onto one line when
it has sufficiently few disjoint lines.  Over prime fields, Sz\H{o}nyi
and Weiner allow a sufficiently small constant multiple of $q^2$
disjoint lines.  This is the scale needed after
our deletion estimate.

The external geometric input is the stability theorem of Sz\H{o}nyi and
Weiner~\cite[Theorem~4]{SzonyiWeiner2012}, also reproduced by H\'eger and
Nagy~\cite[Result~3.6]{HegerNagy2017}.  If $q$ is prime and
$B\subseteq\mathrm{PG}(2,q)$ has
\[
 |B|\le\frac32(q+1)-\beta,
 \qquad
 \Delta<\frac29(\beta+1)^2,
\]
where $\beta>0$ and $\Delta$ counts the projective lines disjoint from
$B$, then some projective line contains at least
$q-2\Delta/(q+1)$ points of $B$.  We use only the following consequence,
with explicit constants to fix its quantifiers.

\begin{lemma}[Prime-plane stability for affine line covers]
\label{lem:prime-stability}
Let $\zeta>0$ and put $\delta_\zeta=\min\{\zeta/4,1/100\}$.
If $q$ is prime and $q\ge\max\{100,4/\zeta\}$, then any set of at most
$q$ distinct affine lines missing fewer than $\delta_\zeta q^2$ affine
points has at least $(1-\zeta)q$ members through a common projective
point.  A point at infinity means that these members are parallel.
\end{lemma}

\begin{proof}
We must account for the points at infinity before applying the
projective statement. Let $h<\delta_\zeta q^2$ be the number of missed
affine points and $d$ the
number of represented directions.  Extending the lines projectively
leaves exactly $q+1-d$ uncovered points at infinity.  Under projective
duality the line set becomes a point set $B$ of size at most $q$, with
\[
 \Delta=h+q+1-d\le h+q+1
\]
skew projective lines.  Set $\beta=\lfloor q/3\rfloor$.  Then
$|B|\le q\le3(q+1)/2-\beta$, and
\[
 \Delta\le h+q+1<\delta_\zeta q^2+q+1
 \le0.0201q^2<\frac{2q^2}{81}
 <\frac29(\beta+1)^2.
\]
Thus the stability theorem applies.  Dualizing its conclusion gives at
least
\[
 q-\frac{2\Delta}{q+1}
 \ge q-\frac{2h}{q+1}-2
 >q-2\delta_\zeta q-2
 \ge(1-\zeta)q
\]
concurrent lines.  The last inequality uses
$2\delta_\zeta\le\zeta/2$ and $2\le\zeta q/2$.
\end{proof}

\begin{corollary}[Concentration of core covers]\label{cor:core-pencil}
Assume the laws of Theorem~\ref{thm:deletion}, and let $q$ tend to infinity
through primes.  For each fixed $\zeta>0$, with probability tending to one,
every set $C\subseteq\mathcal L\setminus D$ of at most $q$ lines covering
$R_D$ has at least $(1-\zeta)q$ members through a common projective point.
\end{corollary}

\begin{proof}
Apply Theorem~\ref{thm:deletion}\textup{(iii)} with
$\delta=\delta_\zeta$, followed by Lemma~\ref{lem:prime-stability}.
\end{proof}
